\section{Réduction de l’universalité de Feigenbaum au certificat hyperbolique de bassin et de transversalité}
\label{sec:p3-final:cierre-feigenbaum}

La famille de Feigenbaum HMT possède un profil analytique, un maximum quadratique,
une dérivée schwarzienne négative et des approximants superstables certifiés jusqu’au niveau
\(8\). Sa promotion à l’universalité exige le certificat
\[
 \ell_u(\partial_t f_t)\ne0
\]
ainsi que l’appartenance au bassin hyperbolique du point fixe. La tangence,
la sortie du bassin ou la présence d’une deuxième valeur propre instable sont
des réfutateurs complets.
